\documentclass[12pt]{article}
\usepackage[margin=1in]{geometry}
\usepackage[colorlinks=true,linkcolor=blue,urlcolor=blue]{hyperref}
\title{Hyperlinks and URLs}
\author{Link Tester}
\begin{document}
\maketitle

\section{Web Links}
Visit the official Python website at \url{https://www.python.org} for the latest
Python releases. The PyMuPDF documentation is available at
\url{https://pymupdf.readthedocs.io}.

\section{Named Links}
For more information, see the \href{https://github.com/pymupdf/PyMuPDF}{PyMuPDF
GitHub repository}. The \href{https://fastapi.tiangolo.com}{FastAPI documentation}
provides excellent tutorials. You can also check the
\href{https://pdfplumber.readthedocs.io}{pdfplumber documentation} for table
extraction details.

\section{Email Links}
Contact the maintainer at \href{mailto:developer@example.com}{developer@example.com}
for support questions.

\section{Internal References}
See Section~\ref{sec:conclusion} for the conclusion.
See Table~\ref{tab:links} for a summary of all URLs.

\begin{table}[h]
\centering
\caption{URL Reference Table}
\label{tab:links}
\begin{tabular}{|l|l|}
\hline
\textbf{Resource} & \textbf{URL} \\
\hline
Python & https://www.python.org \\
PyMuPDF & https://pymupdf.readthedocs.io \\
FastAPI & https://fastapi.tiangolo.com \\
pdfplumber & https://pdfplumber.readthedocs.io \\
\hline
\end{tabular}
\end{table}

\section{Conclusion}
\label{sec:conclusion}
All hyperlinks should be preserved in the Markdown output using the standard
link syntax: [text](url).
\end{document}